\chapter{LITERATURE REVIEW}
\label{chapter2}
This chapter provides a comprehensive review of the literature related to the research domain. The purpose of this chapter is to establish a strong theoretical and conceptual foundation for the study by examining existing research, identifying key contributions, and highlighting gaps in the current knowledge.

A well-structured literature review not only summarizes previous work but also critically evaluates it to develop a coherent understanding of the research problem \cite{hart1998doing,webster2002analyzing}. This chapter also helps justify the need for the present study and supports the formulation of research objectives.

The chapter is organized as follows: Section 2.1 discusses general concepts of literature review and referencing practices. Section 2.2 presents the historical background of the research area. Sections 2.3–2.5 describe major topics relevant to the study. Finally, Section 2.6 summarizes the findings and identifies research gaps.

% -------------------------------------------------
\section{General [Citing and Reference Management]}

The literature review process involves identifying, evaluating, and synthesizing information from relevant sources. It requires critical thinking and the ability to compare different studies to understand their strengths, weaknesses, and contributions. A good literature review should not simply describe previous work but should provide a structured argument that supports the research direction \cite{creswell2014research,booth2016craft}.

Proper citation and referencing are essential components of academic writing. They ensure that original authors receive credit for their work and allow readers to trace the source of information. Common citation styles include IEEE, APA, and Harvard formats. Reference management tools such as Mendeley and EndNote are widely used to organize references and generate bibliographies efficiently \cite{mendeley,endnote}.

Figure \ref{fig:ref_style} illustrates the structure of a typical reference entry.

\begin{figure}[!tbh]
    \centering
    \includegraphics[width=0.7\textwidth]{Images/fig21_referencing.jpg}
    \vspace{0.5em}
    \caption[Structure of a reference entry showing essential components such as author, year, title, and source.]{Structure of a reference entry showing essential components such as author, year, title, and source \cite{apa_author_date}.}
    
    \label{fig:ref_style}
\end{figure}

% -------------------------------------------------
\section{Historical Background}

This section provides an overview of the development of the research area. It includes early studies, significant theoretical advancements, and technological progress. Understanding the historical evolution of a topic helps in identifying trends and recognizing important contributions made by researchers over time.

Research methodology has evolved significantly, transitioning from purely theoretical approaches to more structured and systematic methods \cite{kothari2004research}. Over time, interdisciplinary approaches have also emerged, combining insights from multiple domains to solve complex problems.

% -------------------------------------------------
\section{Topic 1}

This section discusses the first major topic relevant to the research area. It presents various approaches proposed in the literature and compares their effectiveness. Emphasis is given to identifying the advantages and limitations of each method. Table \ref{tab:comparison_methods} summarizes key approaches commonly used in research studies.

\begin{table}[!tbh]
\centering
\caption{Comparison of different approaches in literature}
\label{tab:comparison_methods}
\begin{tabular}{|c|c|c|}
\hline
\textbf{Method} & \textbf{Advantages} & \textbf{Limitations} \\
\hline
Qualitative Research & In-depth analysis & Subjective interpretation \\
Quantitative Research & Objective measurement & Limited flexibility \\
Mixed Methods & Comprehensive insight & Requires more resources \\
\hline
\end{tabular}
\vspace{0.3em}
\end{table}

The choice of approach depends on the nature of the research problem and the objectives of the study \cite{creswell2014research}.

% -------------------------------------------------
\section{Topic 2}

This section focuses on another important aspect of the research. It analyzes different frameworks, models, or methodologies used in similar studies. A comparison of these approaches helps in understanding their applicability and limitations.

Effective research design plays a crucial role in obtaining reliable and valid results. It involves selecting appropriate methods, defining variables, and ensuring that the study is structured logically \cite{kothari2004research}.

% -------------------------------------------------
\section{Topic 3}

This section highlights recent developments and emerging trends in the research domain. With advancements in technology and increasing availability of data, modern research methodologies have become more sophisticated.

Researchers are increasingly adopting data-driven and interdisciplinary approaches to address complex problems. These approaches often provide improved accuracy and efficiency compared to traditional methods \cite{evans2014writing}.

% -------------------------------------------------
\section{Summary and Implications}

This section summarizes the key findings from the literature review and identifies gaps that need further investigation. A good research study should aim to address these gaps and contribute to the existing body of knowledge.

The performance of models and research outcomes is often evaluated using statistical metrics. One commonly used measure is the coefficient of determination ($R^2$), defined as:

\begin{equation}
R^2 = 1 - \frac{\sum_{i=1}^{n} (y_i - \hat{y}_i)^2}
{\sum_{i=1}^{n} (y_i - \bar{y})^2}
\end{equation}

Figure \ref{fig:research_gap} illustrates the concept of identifying research gaps by comparing existing studies and highlighting areas that require further investigation.


\begin{figure}[!tbh]
    \centering
    \includegraphics[width=0.7\textwidth]{Images/fig22_referencing.jpg}
    \vspace{0.5em}
    \caption[Conceptual illustration of research gap showing limitations of existing studies and opportunities for future research.]{Conceptual illustration of research gap showing limitations of existing studies and opportunities for future research \cite{research_gap_blog}.}
    \label{fig:research_gap}
\end{figure}

Based on the review, it is essential to formulate a research problem that addresses the identified limitations. The insights obtained from this chapter provide the foundation for the methodology and analysis presented in the subsequent chapters.